\documentclass[12pt,a4paper]{article}
\usepackage[margin=2.5cm]{geometry}
\usepackage{enumitem}
\setlength{\parindent}{0pt}
\setlength{\parskip}{8pt}
\sloppy
\pagestyle{empty}
\begin{document}
Oleg Tropynin\\
Independent researcher, Karaganda, Kazakhstan\\
tropininoleg@gmail.com \quad ORCID 0009-0006-5431-4285

27 October 2026

The Editor, \textit{Planetary and Space Science}

Transfer of manuscript ICARUS-D-26-00681

Dear Editor,

I submit the manuscript ``Rotation periods and spin senses from Gaia DR3 sparse photometry: tested reliability, 790 candidate new periods and disputed LCDB entries'' for consideration as a research paper.

The manuscript was declined by \textit{Icarus} without review for insufficient novelty. I have revised it to make the new contribution explicit and measurable. Three groups have derived periods and poles from Gaia DR3, and their error rates are known for asteroids with reliable published periods. What a user of a \emph{new} period needs, the error rate for an asteroid without a published period, had not been measured. The paper measures it and delivers what depends on it:
\begin{itemize}[topsep=0pt,itemsep=2pt]
\item 6,238 periods with a probability of being correct, calibrated on half of the LCDB and tested on the held-out half (89\%), against independent ATLAS periods (92.5\%) and on injected lightcurves; 1,680 are absent from the earlier DR3 analyses;
\item spin senses that agree with independent ATLAS poles for 96\%, 1,447 of them for asteroids with no published pole, with their accuracy reported as a function of pole latitude;
\item 39 LCDB periods that disagree with independent survey data, printed in the paper, and 69 more flagged;
\item 790 candidate new periods checked against SsODNet, VizieR and the Minor Planet Bulletin, with an estimated 4--5\% noise and 60--87\% expected to be correct.
\end{itemize}

Since the \textit{Icarus} version I added the independent ATLAS checks of periods and spin senses, the SsODNet novelty search, binomial intervals, the dependence of spin-sense accuracy on pole latitude, a selection explanation of the retrograde excess, and the credit to \v{D}urech \& Hanu\v{s} (2023) for first showing the Yarkovsky segregation in DR3 families, which here serves as a quantified test of the spin senses.

The catalogue, code and validation tables are openly available on Zenodo (doi:10.5281/zenodo.23638281). A short research note on a 600-asteroid pilot (RNAAS, AAS82759) shares only 34 candidate periods with this work. The use of generative AI for code and text is declared in the manuscript.

Thank you for considering this work.

Sincerely,

Oleg Tropynin
\end{document}
